\chapter{Configuración}

\section{Arrancar el orquestador}

El orquestador lee su configuración de variables de entorno:

\begin{tabular}{lp{6.5cm}}
  \textbf{Variable} & \textbf{Contenido} \\
  \texttt{DATABASE\_URL} & Conexión a PostgreSQL. Obligatoria. \\
  \texttt{ORQUESTADOR\_URL} & URL pública que anuncia la AgentCard. Obligatoria. \\
  \texttt{ORQUESTADOR\_TOKENS} & Tokens de los compradores, como \texttt{partyId=token}, separados por comas. \\
  \texttt{ORQUESTADOR\_ORIGENES\_AGENTES} & Orígenes de los que se aceptan AgentCards, separados por comas. Cada uno lleva esquema y host, y el puerto si no es el de por defecto. \\
  \texttt{ORQUESTADOR\_HOST} & Dirección de escucha. Por defecto, \texttt{127.0.0.1}. \\
  \texttt{ORQUESTADOR\_PUERTO} & Puerto de escucha. Por defecto, 8000. \\
\end{tabular}

Un origen se escribe, por ejemplo, como \texttt{https://agentes.ejemplo:8443}. Con las variables definidas:

\begin{verbatim}
uv run python -m orquestador
\end{verbatim}

Si faltan los tokens, nadie puede autenticarse. Si faltan los orígenes, se rechaza cualquier sesión. Un origen o un token mal escrito impide arrancar.

\section{Seguridad del despliegue}

Sin TLS, el token viaja en claro. Por eso el orquestador solo escucha en el propio equipo por defecto. Para exponerlo fuera, pon delante un proxy con TLS.

Los tokens no caducan. Para cambiar uno, actualiza la variable y reinicia el orquestador. Trátalos como contraseñas: guárdalos en \texttt{.env} o en un gestor de secretos, nunca en el repositorio.

\section{Configuración de una sesión}

El coordinador envía una configuración por partida en \texttt{OPEN\_ROUND}:

\begin{tabular}{lp{5cm}l}
  \textbf{Campo} & \textbf{Restricción} & \textbf{Por defecto} \\
  \texttt{fixedTerms} & \texttt{productId}, \texttt{unit} y \texttt{currency} (ISO 4217) & Obligatorio \\
  \texttt{parties} & Dos partes: un \texttt{BUYER} y un \texttt{SELLER} & Obligatorio \\
  \texttt{openingPartyId} & Una de las dos partes & Obligatorio \\
  \texttt{validationMode} & \texttt{STRICT} & Obligatorio \\
  \texttt{maxRetries} & 0 & 2 (se rechaza) \\
  \texttt{turnTimeoutSeconds} & Mayor que 0 & 60 \\
  \texttt{pollIntervalSeconds} & Mayor que 0 y menor que \texttt{turnTimeoutSeconds} & 1 \\
  \texttt{maxTurns} & Al menos 2 & 10 \\
  \texttt{holdTimeoutSeconds} & Mayor que 0 & 1800 \\
\end{tabular}

\section{Límites de esta versión}

\begin{itemize}
  \item \texttt{validationMode = TOLERANT} se rechaza con \texttt{INVALID\_CONFIG}.
  \item Un \texttt{maxRetries} distinto de 0 se rechaza con \texttt{INVALID\_CONFIG}. Como su valor por defecto es 2, hay que enviar 0 de forma explícita.
  \item Un \texttt{OPEN\_ROUND} con \texttt{requestId} se rechaza con \texttt{UNKNOWN\_REQUEST}: solo se admite la primera ronda.
  \item Reenviar una petición con el mismo \texttt{messageId} devuelve un error A2A, en lugar de la Task original. No se escribe nada.
  \item El primer fallo de un agente aborta la sesión: un acto no válido, una Task fallida, un tiempo agotado o un agente que no responde. No hay reintentos.
  \item Las partes aún no reciben \texttt{AGREEMENT\_HELD} ni \texttt{SESSION\_CLOSE}. Una sesión se queda en \texttt{HELD} o en \texttt{CLOSING}, y su resultado solo se ve en la base de datos.
\end{itemize}

\section{Resultado de una sesión}

Mientras no se envía \texttt{SESSION\_CLOSE}, el resultado se lee en la tabla \texttt{sesiones}: la columna \texttt{estado} y, en \texttt{CLOSING}, la columna \texttt{resultado}.

\begin{tabular}{lp{9cm}}
  \textbf{\texttt{outcome}} & \textbf{Qué ha pasado} \\
  \texttt{WITHDRAWN} & Una parte se retiró. \texttt{withdrawnBy} dice cuál. \\
  \texttt{NO\_AGREEMENT} & Se agotaron los turnos sin acuerdo. \\
  \texttt{ABORTED} & Un agente falló. \texttt{lastFailure} dice el primer fallo, con su código y su campo. \\
\end{tabular}

En los tres casos, \texttt{lastOffers} recoge la última oferta de cada parte que llegó a ofertar. Una sesión en \texttt{HELD} tiene un acuerdo provisional: su oferta está en \texttt{oferta\_acordada} y su plazo en \texttt{held\_until}. Los fallos, uno por fila, están en la tabla \texttt{fallos}, y los actos aceptados, en \texttt{actos\_aceptados}.

Ante \texttt{TASK\_INTERRUPTED} con un agente hecho con \texttt{a2a-sdk}, comprueba que el agente publica su Task en \texttt{TASK\_STATE\_SUBMITTED}: una Task sin estado cuenta como interrumpida.

\section{Actualizar una base de datos ya creada}

El orquestador crea las tablas que faltan al arrancar, pero no añade columnas a las que ya existen. Esta versión añade columnas a \texttt{sesiones} y las tablas \texttt{actos\_aceptados}, \texttt{fallos} y \texttt{coerciones}. Una base de desarrollo creada con una versión anterior hay que vaciarla:

\begin{verbatim}
docker compose down -v && docker compose up -d --wait
\end{verbatim}

Esto borra todos sus datos.
